\scanpage{18}
% archive foliation: 17
\[\text{\textcircled{A}}\qquad\boxed{H^{\unsure{*}}(G,A)\Longleftarrow H^p(T/e,\underline{\mathcal H}^{q}(A))}\qquad(T\ne\varnothing)\]
\margin{\unsure{s'apparente} à la suite spectrale de Leray d'un \uline{recouvrement} d'un espace}
où $\underline{\mathcal H}^{q}(A)$ est le \unsure{pré}\uuline{pré}faisceau sur $C_G$ \deleted{défini par $\underline{\mathcal H}^{q}(A)(G/F)=H^q(F,A)$} dont la valeur sur $S\in C_G$ est donnée par
\[
\underline{\mathcal H}^{q}(A)(S)=R_A^q\deleted{\unsure{R}}(\operatorname{Hom}_{G}(S,A))
\]
[ et est donc défini aussi par les conditions de \deleted{\unsure{permuter aux}} transformer sommes \inserted{finies} en produits, et d'être donné, pour $S=F\backslash G$, par $\mathcal H^q(A)(F\backslash G)=H^q(F,A)$ ]

\unsure{Preuve.} \unsure{utiliser} essentiellement les faits suivants

(i) $T\longrightarrow e$ est un morphisme de descente, en particulier \deleted{\unsure{$H^0$}} pour un hyperfaisceau $\Phi$, on a
\[
\Phi(e)=H^0(T/e,\Phi)
\]
(ii) Si $\Phi$ est \inserted{un hyperfaisceau} injectif, alors $H^p(T/e,\Phi)=0$ pour $p>0$ (il suffit de faire $\Phi=\Phi_A$ avec $A=\operatorname{Hom}_{Z}(Z(G),M)=M^G$, alors $\Phi(S)=$ fonc des applications de $S$ dans $M$ ; donc
\[
H^p(T/e,\Phi)=H^p(\text{complexe de \unsure{simplexe} \unsure{simple} défini par }T\ ;\ M)=\unsure{0}
\]
qui est nul si $p>0$.
